\documentclass[12pt]{article}
\usepackage{color,soul}
\usepackage{tikz}
\usetikzlibrary{cd}
% Import preambles and macros for notes
\input{../../../../preambles/notes-preamble.tex}
\input{../../../../preambles/notes-macros.tex}
\input{../../../../preambles/theorem-system-section.tex}

\title{MATH 420 Lecture 21}
\author{Deepak Jassal}
\date{March 30\textsuperscript{th}, 2026}

\begin{document}
\maketitle
\stepcounter{section}
Grothendieck's proof of Riemann-Roch
\begin{enumerate}
    \item Introduced the notion of \underline{schemes} rather than varieties;
    \item Introduced the idea that Riemann-Roch is about \underline{morphisms}, not the actual varieties;
    \item Decomposing morphisms in a certain canonical manner gives ``obvious'' proof.
\end{enumerate}
\propp{Every finitely generated module over a Noetherian ring is Noetherian}{
    Let $M$ be a finitely generated module over a Noetherian ring $R$. If $M$ is generated by a single elements over $R$, say $\alpha$, then $M\cong R/I$ for some ideal $I\subset R$. Since $R$ is Noetherian, so is $R/I$, so we're done. Now induct on teh number $n$ of generators of $M$, say $M=\langle a_1,\dots,a_n \rangle$. Let $N=\langle a_1,\dots,a_{n-1} \rangle$ by the submodule generated by $a_1,\dots,a_{n-1}$. By the IH, $N$ is Noetherian. THen $M/N$ is generated by one element; base case finishes the job. 
}
\prop{Every finitely generated module $M$ over a Noetherian ring $R$ contains a finite chain of submodules $M_r\subset\cdots\subset M_1\subset M_0$ such that $M_i/M_{i+1}$ is isomorphic to $R/P_i$ for some prime ideal $P_i$ or $R$.}
The proof needs the idea of annihilator.
\defn{Annihilator}{
    Let $M$ be an $R$-module. For $x\in M$ the \underline{annihilator} of $x$, denoted ann$(x)$, is the set ann$(x)=\{r\in R:rx=0\}$. 
}
\lemp{}{ann($x$) is an ideal in $R$ ($R$ is commutative).}{
    If $r_1,r_2\in\ann(x)$, then $(r_1+r_2)x=r_1x+r_2x=0$. If $r_1\in\ann(x)$, $r_2\in R$ then $(r_2r_1)x=r_20=0=(r_1r_2)x$.
}
\pf[Proof of Prop 1.2]{
    Let $\mathfrak{a}=\ann(x)$ be maximal among the annihilators of non-zero elements of $M$. We claim that $\ann(x)$ is prime. Suppose $a_1a_2\in\mathfrak{a}$. Then $(a_1a_2)x=0$. Then 
    \[
        \mathfrak{a}\subset(a_1)+\mathfrak{a}\subset \ann(a_2x).
    \]
    If $a_2\in\mathfrak{a}$, then $a_2x\neq0$. Hence, $\mathfrak{a}=\ann(a_2x)$ by maximality. This implies $a_1\in\mathfrak{a}$. To procees, nore that the submodule $Rx$ (orbit of $x$ acted on by $R$) is isomorphic to $R/\ann(x)$. If $M$ is non-zero, then there exists a non-zero $x\in M$ sucj that $\ann(x)$ is maximal, hence prime. Hence $M$ contains a submodule $M_1=Rx$ which is isomorphic to $R/\ann(x)=R/P_1$. Similarily, $M_2$ contains a submodule isomorphic to $R/P_2$. This terminates because $M$ is Noetherian. 
}
\thm{Hilbert Basis Theorem}{
    Every finitely generated algebra over a Noetherian ring $R$ is Noetherian. 
}
\pf[Proof of the Theorem]{
    Let $A$ be a ring, $B$ an $A$-alegbra. We argue by induction over minimum number of generators (for $B$). Note $A[x_1,\dots,x_n]\cong A[x_1,\dots,x_{n-1}][x_n]$. Thus, only the case $n=1$ matters. Let $f(x)=c_rx^r+\cdots+c_1x+c_0\in A[x]$, and $I$ be an ideal of $A{x}$. Let $I(k)$ be the set of elements in $A$ which appear as the leading coefficients of some $f\in I$ with $\deg(f)\leq k$ (can be 0). Then $I(k)$ is an ideal in $A$. Further, $I(k-1)\subset I(k)$, since $c_{k-1}x^{k-1}+\cdots+c_1x+c_0\in I$ implies that $c_{k-1}x^{k}+\cdots+c_1x^2+c_0x\in I$. If $J$ is an ideal in $A[x]$ containing $I$, then $I(k)\subset J(k)$. Further if $I(k)=J(k)$ for all $k\geq0$ then $I=J$. Hence, suppose there exists $f\in J\setminus I$, and $f$ has a minimal degree among all such polynomials. Let $k$ be this degree. Then $I(k)=J(k)$, because for all $g\in J$ of degree at most $k-1$, $g\in I$. If one takes such a $g$ then $f+g\in I$. Then the leading coefficient of $f$, and $f+g$ are the same. Continue on like this\\
    Since $A$ is Noetherian
    \[
        I(1)\subset I(2)\subset \cdots
    \]
    must terminate, say at $I(k)=I(\ell)$ for all $k\geq\ell$. For $i\leq\ell$, there exists a finite generating set $S_i=\{C_{i1},C_{i2},\dots\}$ for $I(i)$. For each $(i,j)$ there exists $f_{ij}\in I$ for degree $i$ and leading coefficient be $C_{ij}$. Let $J=\langle C_{ij} \rangle$. Further, $J(k)=I(k)$ for every $k$, so $I=J$. 
}
\end{document}